\documentclass[11pt,a4paper]{article}
\usepackage[T1]{fontenc}
\usepackage{isabelle,isabellesym}
\usepackage{amsmath,amssymb}
\usepackage{pdfsetup}

\urlstyle{rm}
\isabellestyle{it}

\begin{document}

\title{Verified Compartmental Models: A Reusable Framework\\with SIR Instantiation}
\author{David D.}
\maketitle

\begin{abstract}
We present a reusable Isabelle/HOL framework for certified reasoning about
compartmental ODE models, together with a complete instantiation for the
classical SIR (Susceptible--Infectious--Recovered) epidemic model.

The framework provides a generic conservation theorem for three-compartment
systems, along with monotonicity lemmas derived from the sign of derivatives.
The SIR instantiation proves five key properties: conservation of total
population, monotonicity of the susceptible and recovered compartments,
the epidemic threshold condition, and the stationary infection condition
at the critical susceptible level $S = \gamma/\beta$.

Nonnegativity of compartments is assumed; existence and uniqueness of
solutions are deferred to future work leveraging the AFP ODE infrastructure.
\end{abstract}

\tableofcontents

\section{Introduction}

Compartmental models are a cornerstone of mathematical epidemiology,
pharmacokinetics, and systems biology. The SIR model, introduced by
Kermack and McKendrick~\cite{kermack1927}, partitions a population into susceptible~(S),
infectious~(I), and recovered~(R) compartments governed by a system of ODEs.

This development formalizes the key qualitative properties of the SIR model
in Isabelle/HOL, structured as a reusable compartmental model framework
with the SIR model as its first instantiation.

\paragraph{Contributions.}
\begin{itemize}
\item A generic three-compartment conservation theorem.
\item Monotonicity corollaries for compartments with signed derivatives.
\item A locale-based SIR formalization proving conservation, monotonicity,
      the epidemic threshold, and the stationary infection condition.
\end{itemize}

% generated text of all theories
\input{session}

\bibliographystyle{abbrv}
\bibliography{root}

\end{document}
